\documentclass[12pt]{article}
\usepackage[utf8]{inputenc}
\usepackage[T1]{fontenc}
\usepackage{hyperref}
\title{Urningangg Language Resources for K-12 Curriculum Development}
\author{Compiled by Claude Sonnet 5 (claude-sonnet-5) for Andr\'as Kornai}
\date{}
\begin{document}
\maketitle
\section{Introduction}
Urningangg (also Urningangk, Urningangga; ISO 639-3 urc; Glottocode urni1239) is an extinct language of western Arnhem Land, Northern Territory, Australia, belonging to the Giimbiyu group together with Mangerr (Mengerrdji) and Erre. No native speakers remain and the search found no school programme, primer or revitalization curriculum. Existing material is documentary and academic only.
\section{Primary Educational Resources}
No primers, textbooks or school materials were found in search results.
\section{Linguistic and Reference Materials}
\begin{itemize}
\item Giimbiyu language, Wikipedia overview of the group: \url{https://en.wikipedia.org/wiki/Giimbiyu_language}
\item Glottolog entry for Urningangg: \url{https://glottolog.org/resource/languoid/id/urni1239}
\item Grambank entry for Urningangg: \url{https://grambank.clld.org/languages/urni1239}
\item Lauren Campbell (2006), honours thesis, sketch grammar of Urningangk, Erre and Mengerrdji (reference record): \url{https://glottolog.org/resource/reference/id/32813}
\end{itemize}
Search summaries also mention Bruce Birch's 2006 dictionary of Erre, Mengerrdji and Urningangk, but no URL for it was seen.
\section{Religious and Community Resources}
None found.
\section{Digital and Audio Resources}
None found. A Western Arnhem Land cultural PDF on the Urningangga was listed: \url{https://hgcosmos.org/assets/culture/AZ/U/Urningangga.pdf}
\section{Curriculum Development Context}
Any curriculum would have to be reconstructive and built from the academic descriptions above, ideally with the descendant community and regional Arnhem Land language bodies. The materials found are thin and not classroom-ready.
\end{document}